\documentclass{article}

\PassOptionsToPackage{numbers, sort, compress}{natbib}
\usepackage[main,final]{neurips_2025}

\usepackage[T2A]{fontenc}
\usepackage[utf8]{inputenc}
\usepackage[english,russian]{babel}
\usepackage{hyperref}
\usepackage{url}
\usepackage{booktabs}
\usepackage{nicefrac}
\usepackage{microtype}
\usepackage{xcolor}

%%%

\usepackage{subcaption}
\usepackage{graphicx}
\usepackage{multirow}
\usepackage{amsmath,amssymb,amsfonts}
\usepackage{amsthm}
\usepackage{mathrsfs}
\usepackage{xcolor}
\usepackage{textcomp}
\usepackage{manyfoot}
\usepackage{algorithm}
\usepackage{algorithmicx}
\usepackage{algpseudocode}
\usepackage{listings}

\newtheorem{theorem}{Теорема} % continuous numbers
%%\newtheorem{theorem}{Theorem}[section] % sectionwise numbers
%% optional argument [theorem] produces theorem numbering sequence instead of independent numbers for Proposition
\newtheorem{proposition}[theorem]{Утверждение}%
\newtheorem{lemma}{Лемма}%
%%\newtheorem{proposition}{Proposition} % to get separate numbers for theorem and proposition etc.

\newtheorem{example}{Пример}
\newtheorem{remark}{Замечание}

\newtheorem{definition}{Определение}
\newtheorem{assumption}{Предположение}

% neurips_2025.sty печатает заголовок аннотации по-английски; переопределяем на русский.
\renewenvironment{abstract}%
{%
  \vskip 0.075in%
  \centerline%
  {\large\bf Аннотация}%
  \vspace{0.5ex}%
  \begin{quote}%
}
{
  \par%
  \end{quote}%
  \vskip 1ex%
}

%%%


\title{CRISP: самоигра Challenger--Reasoner на основе корпуса документов для дообучения больших языковых моделей в предметных областях}


% The \author macro works with any number of authors. There are two commands
% used to separate the names and addresses of multiple authors: \And and \AND.
%
% Using \And between authors leaves it to LaTeX to determine where to break the
% lines. Using \AND forces a line break at that point. So, if LaTeX puts 3 of 4
% authors names on the first line, and the last on the second line, try using
% \AND instead of \And before the third author name.


\author{
  Роман Акрамов\\
  МФТИ\\
  Москва, Россия\\
  \texttt{akramovromanr@gmail.com}\\
  \And
  Роман Валерьевич Ищенок\\
  Научный руководитель\\
  МФТИ\\
  Москва, Россия\\
}


\begin{document}


\maketitle

\begin{abstract}
    TODO
\end{abstract}

\textbf{Ключевые слова:} самоигра, обучение с подкреплением, большие языковые модели, курирование данных, GRPO, синтез данных на основе корпуса

\section{Введение}\label{sec:intro}

Дообучение (post-training) больших языковых моделей (LLM) -- обучение с учителем (SFT) на отобранных демонстрациях и обучение с подкреплением с проверяемыми наградами (RLVR) с помощью GRPO~\cite{shao2024deepseekmath} и Dr.\,GRPO~\cite{liu2025drgrpo} -- хорошо отработано для узкого набора изученных областей: школьной математики, генерации кода и диалоговых ассистентов. Однако эти подходы трудно применить в областях, где нет готового набора пар <<вопрос--ответ>> (QA) с надёжными эталонными ответами и откалиброванной сложностью. Узким местом стали уже не алгоритмы обучения, а получение данных, тогда как ручная разметка экспертами остаётся слишком дорогой.

У существующих автоматических подходов есть известные ограничения. Дистилляция из модели-учителя~\cite{wang2022selfinstruct,xu2023wizardlm,luo2023wizardmath} ограничена возможностями учителя и его знанием предметной области и часто неэффективна в узких доменах. Сбор сырого корпуса~\cite{penedo2024fineweb} даёт данные для предобучения, но не пары <<вопрос--ответ>>. Генерация вопросов по корпусу обычно не калибрует их сложность, а самоигра без опоры на документы~\cite{huang2025rzero,zhao2025absolute} подвержена галлюцинациям и дрейфу распределения.

Самоигра \emph{на основе корпуса}, в которой генератор вопросов работает с реальным документом, а эталонные ответы извлекаются из текста, сочетает фактическую обоснованность с адаптивной сложностью: генератор вопросов получает награду в зависимости от дисперсии успешности отвечающей модели. Эту идею предложили авторы SPICE~\cite{spice} в виде алгоритма обучения, но код и данные не опубликовали. Мы заново реализуем и расширяем этот подход: его схема награды и взаимодействие двух ролей становятся основой более широкой сквозной <<фабрики данных>>, которую мы называем \textbf{CRISP} (\textbf{C}orpus-grounded Challenger--\textbf{R}easoner Self-Play). Вместо одного обученного чекпойнта CRISP превращает набор исходных PDF-документов предметной области в два переиспользуемых артефакта: корпус для предобучения и QA-датасет, пригодный без изменений и для SFT, и для RL.

\textbf{Описание проекта.} Конвейер состоит из четырёх этапов:
\begin{enumerate}
    \item \textbf{OCR.} Исходные PDF преобразуются в структурированный Markdown с сохранением математической нотации, таблиц и иерархии документа.
    \item \textbf{Очистка, разбиение на фрагменты и фильтрация.} Удаляются артефакты OCR и издательские метаданные, текст разбивается на смысловые фрагменты (чанки), каждый фрагмент оценивается по качеству и размечается темами с помощью LLM. В результате получается чистый корпус для предобучения, который также служит входом для самоигры.
    \item \textbf{Самоигра на основе корпуса.} Одна LLM обучается в двух ролях с отдельными LoRA-адаптерами: \emph{Challenger} читает фрагмент и формулирует вопрос с эталонным ответом, опирающимся на текст, а \emph{Reasoner} отвечает на него без доступа к документу. Challenger получает награду, зависящую от дисперсии: она максимальна, когда Reasoner решает вопрос примерно в половине попыток, поэтому сгенерированные вопросы держатся на границе возможностей модели.
    \item \textbf{Курирование.} Журнал вопросов самоигры заново оценивается итоговой моделью, фильтруется по сложности и проходит через LLM-судью с несколькими критериями (отдельную, более сильную модель, а не саму обучаемую). На выходе получается готовый к обучению QA-датасет.
\end{enumerate}
Конвейер применяется к трём русскоязычным STEM-областям: общим естественно-научным и математическим дисциплинам университетского уровня, линейной алгебре, теории вероятностей и математической статистике. Корпус для предобучения оценивается продолженным предобучением в сравнении с открытыми датасетами для предобучения, а QA-датасет -- дообучением с помощью GRPO в сравнении с открытыми эталонными датасетами.

\textbf{Вклад.} Основные результаты работы:
\begin{itemize}
    \item \textbf{Конвейер подготовки данных для предобучения}, который принимает PDF-документы предметной области и строит корпус с оценками качества и тематической разметкой для продолженного предобучения.
    \item \textbf{Алгоритм самоигры}, обучающий одну LLM в двух ролях (Challenger и Reasoner) с помощью Dr.\,GRPO и награды, зависящей от дисперсии, и генерирующий вопросы на границе возможностей модели.
    \item \textbf{Конвейер курирования данных для выравнивания}, который превращает журнал самоигры в единый QA-датасет, пригодный без изменений и для SFT, и для RL, с помощью многокритериальной фильтрации LLM-судьёй.
\end{itemize}

\textbf{Структура работы.} В разделе~\ref{sec:rw} приведён обзор литературы, в разделе~\ref{sec:prelim} описаны GRPO и Dr.\,GRPO, в разделе~\ref{sec:method} -- предлагаемый конвейер, в разделе~\ref{sec:exp} -- эксперименты, в разделах~\ref{sec:disc}--\ref{sec:concl} -- обсуждение результатов и заключение.

\section{Обзор литературы}\label{sec:rw}

\textbf{Тема \#1.}
TODO

\textbf{Тема \#2.}
TODO

\section{Предварительные сведения}\label{sec:prelim}

\subsection{Обозначения}

TODO

\subsection{Предположения}

TODO

\section{Метод}\label{sec:method}

TODO

\section{Эксперименты}\label{sec:exp}

TODO

\section{Обсуждение}\label{sec:disc}

TODO

\section{Заключение}\label{sec:concl}

TODO


%%%%%%%%%%%%%%%%%%%%%%%%%%%%%%%%%%%%%%%%%%%%%%%%%%%%%%%%%%%%

\bibliographystyle{unsrtnat}
\bibliography{references}

%%%%%%%%%%%%%%%%%%%%%%%%%%%%%%%%%%%%%%%%%%%%%%%%%%%%%%%%%%%%

\newpage
\appendix
\section{Приложение}\label{app}

\subsection{Дополнительные эксперименты}\label{app:exp}

TODO

\end{document}
